\section{FreeGrasp：机器人抓取中的 Visual Prompting}

\subsection{问题设定}

FreeGrasp 解决的是机器人自由形式抓取（free-form grasping）问题：用户用自然语言给出模糊的抓取指令（如"抓底部的红色玩具车"），机器人需要理解指令、定位目标物体、处理遮挡，最终完成抓取。

\begin{knowledgebox}{遮挡处理是关键挑战}
在真实抓取场景中，目标物体常被其他物体遮挡。FreeGrasp 设计了一个完整的 Pipeline 来处理遮挡：先识别目标→判断是否被遮挡→移开遮挡物→再抓取目标。这个过程需要 VLM 具备多步推理能力。
\end{knowledgebox}

\subsection{Pipeline 设计}

FreeGrasp 的处理流程：

\begin{enumerate}
    \item \textbf{输入}：RGBD 图像（RGB 用于视觉推理，Depth 用于机械臂执行）+ 自然语言指令
    \item \textbf{全图检测}：用预训练的 CV 模型检测全图所有物体，获取每个物体的中心点
    \item \textbf{Visual Marking}：在中心点位置标注数字编号
    \item \textbf{VLM 推理}：使用结构化 Prompt，让 VLM 基于带编号的图片进行多步推理：
    \begin{itemize}
        \item 识别目标物体（哪个编号是目标？）
        \item 判断是否被遮挡（哪些物体遮挡了目标？）
        \item 规划移除遮挡物的顺序
    \end{itemize}
    \item \textbf{分割与抓取}：对目标物体做分割，传统算法估计抓取姿态，机械臂执行
\end{enumerate}

\begin{importantbox}{Free-form 的含义}
FreeGrasp 中"Free"指的是用户输入的自由度——可以用任意自然语言描述抓取目标，无需结构化的指令格式。这得益于 VLM 强大的自然语言理解能力。
\end{importantbox}

\subsection{本章小结}

FreeGrasp 将 Visual Prompting 技术应用于机器人抓取领域，通过检测→编号→VLM 推理→分割→执行的 Pipeline，实现了基于自由形式自然语言指令的物体抓取，其中遮挡检测与处理是核心创新。


